% AUTHORED
\section{Backpropagation / autograd over a graph}\label{sec:05}
\conceptbadge{05}

Backpropagation computes derivatives of a scalar loss with respect to every
parameter that influenced it. The cleanest way to see it is as a computation
graph. Each node stores a value from the forward pass, a gradient accumulated
during the backward pass, and a small function that knows how to pass gradients
to its parents.

\subsection*{Local derivative times upstream gradient}
Suppose a node computes $z=f(x,y)$ and the final loss is $L$. During backward,
the node receives an upstream gradient $\partial L/\partial z$. It sends to $x$
the product
\[
  \frac{\partial L}{\partial x}
  =
  \frac{\partial L}{\partial z}
  \frac{\partial z}{\partial x},
\]
and similarly for $y$. This is the chain rule. The operation only needs its
local derivative; the upstream gradient carries everything that happened later.

\subsection*{Two basic backward rules}
For addition,
\[
  z = x + y,
  \qquad
  \frac{\partial z}{\partial x}=1,
  \quad
  \frac{\partial z}{\partial y}=1.
\]
So the same upstream gradient flows to both inputs. For multiplication,
\[
  z = xy,
  \qquad
  \frac{\partial z}{\partial x}=y,
  \quad
  \frac{\partial z}{\partial y}=x.
\]
The upstream gradient is scaled by the other input's forward value.

\subsection*{Accumulation}
A value can feed into more than one later operation. Its final gradient is the
sum of all paths by which it affects the loss. Autograd engines therefore add
to gradients rather than overwriting them. This is also why training loops need
a clear ``zero gradients'' step before each new backward pass.

\subsection*{Reverse topological order}
Backward must run after all downstream gradients for a node are known. A
topological ordering lists parents before children in the forward graph; reverse
topological order visits children before parents during backward. P1 asks you to
build exactly this: trace the graph from the loss, seed
$\partial L/\partial L = 1$, then call each node's backward function in reverse
topological order.

\subsection*{From scalars to tensors}
The project begins with scalar-style autograd because it makes the chain rule
visible. Later primitives can use arrays, but the principle is unchanged:
forward operations record enough context to compute local derivatives, and the
backward pass pushes gradients toward parameters.

\medskip
\noindent\textit{Go deeper:} Karpathy's \texttt{micrograd}
(\texttt{https://github.com/karpathy/micrograd}) is a minimal scalar autograd
engine that implements exactly this reverse-mode graph; Goodfellow, Bengio \&
Courville, \emph{Deep Learning}, Ch.~6.5
(\texttt{https://www.deeplearningbook.org}) presents back-propagation more
formally.
